\chapter*{Förord}
\addcontentsline{toc}{chapter}{Förord}
\markboth{Förord}{Förord}

Den här boken handlar om två sätt att lösa samma problem.

Första delen löser det i mjukvara. Du bestämmer själv var en frame börjar, hur lång den är, hur den
skyddas mot fel, och vad som ska hända när ett svar uteblir. Varje mekanism blir din egen, och
varje mekanism blir kod du får skriva och underhålla. Det är lärorikt på ett sätt som är svårt att
ersätta: man förstår vad ett längdfält är till för först när man glömt det.

Andra delen kastar allt det och låter hårdvaran göra jobbet. CAN är från 1986, det sitter i varje
bil byggd sedan mitten av nittiotalet, och det löser nästan exakt de problem du just löst själv
--- fast i kisel, och bättre. Det som återstår för mjukvaran är att prata med kontrollern.

Tredje delen kopplar ihop det med verkligheten: en mikrokontroller, ett FPGA, fyra ledningar
mellan dem, och en buss som en bussanalysator kan läsa.

\section*{Det som gör kursen ovanlig}

CAN-kontrollern du skriver drivrutin mot finns inte färdig. Den konstrueras i VHDL av en
parallellklass, samtidigt som du skriver mjukvaran mot den.

Ingen av klasserna kan alltså kompilera mot den andras kod. Det enda som binder ihop de två
halvorna är två dokument --- en registerkarta och ett SPI-protokoll --- och de är projektets
egentliga kravspecifikation.

\begin{aside}
\note Att arbeta mot ett kontrakt i stället för mot en implementation är inte en kursövning. Det
är hur inbyggd utveckling fungerar när hårdvara och mjukvara byggs parallellt, vilket är det
normala fallet och inte undantaget. Det mesta av det som är svårt i den här kursen är svårt av den
anledningen.
\end{aside}

\section*{Bokens upplägg}

Ett kapitel per föreläsning, femton stycken, i tre delar:

\begin{booktable}{@{}lL@{}}
\thead{Kapitel} & \thead{Del} \\ \midrule
1--3 & \textbf{Ett eget protokoll.} Frames, byte-parsning, bussar, routing och tillförlitlighet. \\
4--12 & \textbf{CAN-drivrutinen.} Registerkartan, arkitekturen, interfacet, stubben, SPI,
  transportsömmen, drivrutinen, testningen och hårdvaran. \\
13--15 & \textbf{Hårdvaran.} Bring-up, bussanalys, och repetition inför duggan. \\
\end{booktable}

Varje kapitel öppnar med vad det innehåller, avslutas med en sammanfattning, och har därefter
sina övningar. \textbf{Lösningarna till varje övning som inte ber om kod står i appendix~A} ---
alla resonemangs-, räkne- och avkodningsuppgifter, och de delar av kodövningarna som är resonemang.
Koden som kodövningarna ber om besvaras inte där: för den är den utdelade testsviten svaret, och
koden publiceras i kursrepot efter respektive föreläsning.

Övningsduggan med fullständiga svar ligger sist, i appendix.

\section*{Förkunskaper}

Boken förutsätter klasser, arv, virtuella funktioner, abstrakta basklasser och factory-mönstret.
Den repeterar dem inte. Den förutsätter också grundläggande vana vid Git, en terminal och
\code{make}.

Inget om CAN, om hårdvaruregister eller om testning förutsätts.

\section*{Ytterligare litteratur}

En bok till hör till kursen, och den behövs inte för att klara den: \textit{CAN --- bussen, framen
och kontrollern}, kursens egen bok om CAN som protokoll. Den tar bussen, framen, bitstoppningen,
CRC-15, arbitreringen, bittajmingen och felhanteringen, med övningar och svar, och är att betrakta
som det här kapitel 4:s långa version. Den ligger i ett eget repo, eftersom den delas med
hårdvarukursen och därför inte nämner någon kurs.

